\documentclass[10pt,a4paper]{amsart}
\usepackage[margin=19mm]{geometry}
\usepackage{amsmath,amssymb,mathtools}
\usepackage[scr=rsfs,cal=euler]{mathalfa}
\usepackage{xfrac}
\usepackage[unicode,hidelinks,bookmarks=false]{hyperref}
\newcommand{\ie}{i.e.\ }
\newcommand{\cf}{cf.\ }
\newcommand{\resp}{respectively\ }
\newcommand{\Subsection}{\subsection}
\providecommand{\nobreakdash}{\nobreak\hbox{-}}
\setlength{\emergencystretch}{2em}
\multlinegap0pt
\usepackage{amssymb}
\usepackage{mathtools}
\usepackage{float}
\usepackage[shortlabels]{enumitem}
\usepackage{xcolor}
\newtheorem{lemma}{Lemma}
\title{Diffusion Computation versus Quantum Computation: A Comparative Model for Order Finding and Factoring}
\author{Carlos A. Cadavid, Paulina Hoyos, Jay Jorgenson, Lejla Smajlovi\'c, J. D. V\'elez}
\date{Source: arXiv:2601.02518v1 (CC BY 4.0)}
\begin{document}
\maketitle

\noindent\emph{Source and licence.} Diffusion Computation versus Quantum Computation: A Comparative Model for Order Finding and Factoring. Version arXiv:2601.02518v1 (CC BY 4.0), distributed by arXiv under the Creative Commons Attribution 4.0 International licence, http://creativecommons.org/licenses/by/4.0/, as recorded in the arXiv licence metadata for this exact version and shown on the abstract page https://arxiv.org/abs/2601.02518v1. Full licence notice: \texttt{licenses/arxiv-2601.02518-CC-BY-4.0.txt}.\par\smallskip

\noindent\textit{Editorial scope.} Counted unit: the article's lemma lem:IE-mobius-finite-primes (source0.tex lines 1195--1219). The article's complete original proof of it is delivered with it (source0.tex lines 1221--1296, one \texttt{proof} environment of 76 lines). No mathematics is written, repaired or re-derived: the statement and the proof are the article's own lines, in the article's own notation, and no retained line is substituted or re-flowed.  No further source-local context is needed: every cross-reference of every retained line resolves inside the two delivered spans.  Nothing was omitted from any retained span, so the package carries no omission record.  No retained line contains a citation, so the wrapper carries no bibliography and no bibliography entry.  No retained line contains a figure environment, an image-inclusion command or drawing code, so no image is delivered and the package declares no asset dependency.  Editorial formatting only, in addition: an AMS \texttt{amsart} preamble that restates the article's own declarations and macros used by the retained lines, namely the environments \texttt{lemma} on separate counters, together with the standard packages those lines need.  The retained environments print in this document as Lemma 1 in order of appearance, whereas the article prints the same items with the numbering of its own full document.  The complete native source of the article as distributed in the arXiv e-print for this exact version is bundled unchanged as \texttt{sources/192155/source0.tex}.
\begin{lemma}
\label{lem:IE-mobius-finite-primes}
Let $\mathcal{P}$ be a fixed finite set of primes, and let $U_1,\dots,U_s$ be independent random
variables uniformly distributed on $\Omega=\{1,2,\dots,Q\}$. For each $\ell\in\mathcal{P}$ let
\[
E_\ell = \{\ell \mid U_1\cdots U_s\},
\qquad\text{so that}\qquad
E_\ell^{\,c}=\{\ell \nmid U_1\cdots U_s\}.
\]
Then inclusion--exclusion can be written as
\[
\Pr\!\left(\bigcap_{\ell\in\mathcal{P}} E_\ell^{\,c}\right)
=\sum_{d\mid \prod_{\ell\in\mathcal{P}}\ell} \mu(d)\,\Pr(d\mid U_1,\dots,d\mid U_s),
\]
where $\mu$ is the M\"obius function. Moreover,
\[
\Pr(d\mid U_1,\dots,d\mid U_s)=\Bigl(\frac{\lfloor Q/d\rfloor}{Q}\Bigr)^s,
\]
and therefore
\[
\lim_{Q\to\infty}\Pr\!\left(\bigcap_{\ell\in\mathcal{P}} E_\ell^{\,c}\right)
=\sum_{d\mid \prod_{\ell\in\mathcal{P}}\ell}\frac{\mu(d)}{d^s}
=\prod_{\ell\in\mathcal{P}}\left(1-\frac{1}{\ell^s}\right).
\]
\end{lemma}

\begin{proof}
For each $\ell\in\mathcal{P}$ the event $E_\ell$ occurs iff $\ell$ divides at least one of the
$U_i$, hence by De Morgan,
\[
\bigcap_{\ell\in\mathcal{P}}E_\ell^{\,c}
=\Bigl(\bigcup_{\ell\in\mathcal{P}}E_\ell\Bigr)^{c}.
\]
By De Morgan's law,
\[
\bigcap_{\ell\in\mathcal{P}}E_\ell^{\,c}
=\Bigl(\bigcup_{\ell\in\mathcal{P}}E_\ell\Bigr)^{c},
\]
hence
\[
\Pr\!\left(\bigcap_{\ell\in\mathcal{P}}E_\ell^{\,c}\right)
=1-\Pr\!\left(\bigcup_{\ell\in\mathcal{P}}E_\ell\right).
\]
Applying inclusion--exclusion to the union gives
\[
\Pr\!\left(\bigcup_{\ell\in\mathcal{P}}E_\ell\right)
=\sum_{\varnothing\neq S\subseteq \mathcal{P}}(-1)^{|S|+1}
\Pr\!\left(\bigcap_{\ell\in S} E_\ell\right).
\]
Substituting this into the previous line and simplifying the sign yields
\[
\Pr\!\left(\bigcap_{\ell\in\mathcal{P}}E_\ell^{\,c}\right)
=1+\sum_{\varnothing\neq S\subseteq \mathcal{P}}(-1)^{|S|}
\Pr\!\left(\bigcap_{\ell\in S} E_\ell\right).
\]
Finally, we absorb the term \(1\) into the sum by allowing \(S=\varnothing\):
since \(\bigcap_{\ell\in\varnothing}E_\ell=\Omega\) and \(\Pr(\Omega)=1\), we obtain
\[
\Pr\!\left(\bigcap_{\ell\in\mathcal{P}}E_\ell^{\,c}\right)
=\sum_{S\subseteq \mathcal{P}}(-1)^{|S|}
\Pr\!\left(\bigcap_{\ell\in S} E_\ell\right),
\]
where the term \(S=\varnothing\) equals \(1\).

Fix a subset $S\subseteq\mathcal{P}$ and put $d_S=\prod_{\ell\in S}\ell$ (squarefree). Since
$E_\ell=\{\ell\mid U_1\cdots U_s\}$ and the primes in $S$ are distinct, we have
\[
\bigcap_{\ell\in S} E_\ell
=\{d_S \mid U_1\cdots U_s\}
=\bigcup_{\substack{i_1,\dots,i_s\in\{0,1\}\\ i_1+\cdots+i_s\ge 1}}
\Bigl\{\prod_{\ell\in S}\ell^{\,i_1+\cdots+i_s}\mid U_1\cdots U_s\Bigr\},
\]
and in particular (because $d_S$ is squarefree) the condition $d_S\mid U_1\cdots U_s$ is
equivalent to requiring that \emph{each} prime $\ell\in S$ divides at least one of the $U_i$.
Re-indexing the inclusion--exclusion sum by squarefree divisors $d$ of
$D=\prod_{\ell\in\mathcal{P}}\ell$ gives exactly
\[
\Pr\!\left(\bigcap_{\ell\in\mathcal{P}} E_\ell^{\,c}\right)
=\sum_{d\mid D}\mu(d)\,\Pr(d\mid U_1,\dots,d\mid U_s),
\]
since for squarefree $d$ with prime set $S$ one has $\mu(d)=(-1)^{|S|}$, and $\mu(d)=0$ for
non-squarefree $d$.

Finally, for a fixed integer $d\ge 1$,
\[
\Pr(d\mid U_i)=\frac{\#\{1\le n\le Q: d\mid n\}}{Q}=\frac{\lfloor Q/d\rfloor}{Q},
\]
and by independence,
\[
\Pr(d\mid U_1,\dots,d\mid U_s)=\Bigl(\frac{\lfloor Q/d\rfloor}{Q}\Bigr)^s.
\]
Letting $Q\to\infty$ yields $\lfloor Q/d\rfloor/Q\to 1/d$, hence
\[
\lim_{Q\to\infty}\Pr\!\left(\bigcap_{\ell\in\mathcal{P}} E_\ell^{\,c}\right)
=\sum_{d\mid D}\frac{\mu(d)}{d^s}.
\]
The last equality
\(
\sum_{d\mid D}\mu(d)d^{-s}=\prod_{\ell\in\mathcal{P}}(1-\ell^{-s})
\)
is the standard multiplicative factorization over squarefree divisors of $D$.
\end{proof}

\end{document}
